\documentclass[12pt]{article}
\usepackage[utf8]{inputenc}
\usepackage[T1]{fontenc}
\usepackage[margin=1in]{geometry}
\usepackage{setspace}
\usepackage{xcolor}
\usepackage[colorlinks=true, urlcolor=blue!60!black, linkcolor=blue!60!black]{hyperref}
\setlength{\parindent}{2em}
\setlength{\parskip}{1.25em}
\renewcommand{\baselinestretch}{1.0}
\begin{document}
\begin{center}
{\LARGE \textit{Cedar Point Nursery v. Hassid}, 594 U.S. 139 (2021)\par}
\vspace{0.8em}
{\large Case Brief\par}
\end{center}
\vspace{1em}

\section{Memory Jogger}
A California rule letting union organizers enter farms for up to three hours a day, 120 days a year, was held 6--3 to be a \emph{per se} physical taking: a regulation that appropriates the owner's right to exclude for third parties requires just compensation, however intermittent the access, subject to limited exceptions (\href{https://supreme.justia.com/cases/federal/us/594/20-107/}{Justia}).

\section{Detailed Case Facts}
Under California's Agricultural Labor Relations Act of 1975, the Agricultural Labor Relations Board (ALRB) promulgated a regulation granting labor organizations a ``right to take access'' onto an agricultural employer's private premises to meet with employees and solicit union support (Cal.\ Code Regs., tit.~8, \S~20900(e)) (\href{https://www.law.cornell.edu/supremecourt/text/20-107}{Cornell LII}). Upon filing written notice with the Board and serving the employer, union organizers could enter the employer's property for up to three hours per day---one hour before work, one hour during lunch, and one hour after work---during up to four 30-day periods (120 days) per calendar year. Interfering with access constituted an unfair labor practice carrying state sanctions.

Two California agricultural employers whose workers lived entirely off-premises challenged the regulation (\href{https://www.law.cornell.edu/supremecourt/text/20-107}{Cornell LII}):
\begin{itemize}
    \item \textbf{Cedar Point Nursery:} A northern California strawberry grower employing over 400 seasonal and roughly 100 full-time workers. In October 2015, United Farm Workers (UFW) organizers entered Cedar Point's trim shed at 5:00 a.m.\ without prior notice, using bullhorns and disrupting operations; both sides filed unfair labor practice charges.
    \item \textbf{Fowler Packing Company:} A Fresno-based table-grape and citrus grower and shipper with 1,800--2,500 field workers and roughly 500 packing-facility employees. In July 2015, Fowler blocked UFW organizers from entering its property, prompting an unfair labor practice charge that the union later withdrew.
\end{itemize}

\section{Procedural History}
Cedar Point and Fowler sued ALRB officials in the U.S. District Court for the Eastern District of California, seeking declaratory and injunctive relief on the ground that the access regulation appropriated an uncompensated easement, effecting a \emph{per se} physical taking under the Fifth and Fourteenth Amendments. The District Court denied a preliminary injunction and dismissed the complaint, holding that because access was neither permanent nor continuous, the claim had to be evaluated under the multifactor balancing test of \textit{Penn Central Transportation Co. v. New York City}, 438 U.S. 104 (1978), which the growers did not plead. A divided Ninth Circuit panel affirmed (923 F.3d 524 (2019)), reasoning that \emph{per se} treatment applies only to permanent physical occupations allowing access ``24 hours a day, 365 days a year.'' Rehearing en banc was denied over a dissent by Judge Ikuta joined by seven judges (956 F.3d 1162 (2020)), and the Supreme Court granted certiorari (\href{https://supreme.justia.com/cases/federal/us/594/20-107/}{Justia}).

\section{Judicial Votes}
The Supreme Court \textbf{reversed and remanded} the Ninth Circuit's decision (923 F.3d 524) by a \textbf{6--3} vote. The case was argued on March 22, 2021, and decided on June 23, 2021 (\href{https://supreme.justia.com/cases/federal/us/594/20-107/}{Justia}; \href{https://www.law.cornell.edu/supremecourt/text/20-107}{Cornell LII}):

\begin{itemize}
    \item \textbf{Majority Opinion:} Chief Justice Roberts delivered the opinion of the Court, joined by Justices Thomas, Alito, Gorsuch, Kavanaugh, and Barrett.
    \item \textbf{Concurring Opinion:} Justice Kavanaugh filed a concurring opinion.
    \item \textbf{Dissenting Opinion:} Justice Breyer filed a dissenting opinion, joined by Justices Sotomayor and Kagan.
\end{itemize}

\section{Holding}
\textbf{Questions Presented:}

Does a state regulation granting union organizers a right to physically enter an agricultural employer's private property for up to three hours per day, 120 days per year, constitute a \emph{per se} physical taking under the Fifth and Fourteenth Amendments, or is it a use restriction evaluated under \textit{Penn Central} balancing?

\textbf{Holding:} \textbf{Per Se Physical Taking:} California's access regulation appropriates a right to physically invade the growers' property---transferring to third parties the owners' fundamental right to exclude---and therefore constitutes a \emph{per se} physical taking requiring just compensation under the Fifth and Fourteenth Amendments, regardless of whether the authorized access is intermittent or temporary (\href{https://www.law.cornell.edu/supremecourt/text/20-107}{Cornell LII}).

\section{Analysis of Opinions}
\subsection{Majority Opinion (Chief Justice Roberts)}

\subsection{Physical Appropriations vs.\ Use Restrictions}
Writing for a 6--3 majority, Chief Justice Roberts drew a sharp structural divide in Takings Clause doctrine. When government restricts how an owner may use her own property (e.g., zoning or commercial limits), courts apply the flexible balancing test of \textit{Penn Central} to determine whether the regulation goes ``too far'' (\textit{Pennsylvania Coal Co. v. Mahon}, 260 U.S. 393, 415 (1922)). By contrast, when government physically appropriates property for itself or a third party---whether by eminent domain, physical seizure (\textit{United States v. Pewee Coal Co.}, 341 U.S. 114 (1951)), or administrative regulation (\textit{Horne v. Department of Agriculture}, 576 U.S. 351 (2015))---a categorical \emph{per se} rule applies: the government must pay for what it takes, and \textit{Penn Central} has no place (\href{https://www.law.cornell.edu/supremecourt/text/20-107}{Cornell LII}).

\subsection{Appropriating the Right to Exclude}
Because the right to exclude is ``one of the most treasured'' and ``essential sticks in the bundle of rights'' (\textit{Loretto v. Teleprompter Manhattan CATV Corp.}, 458 U.S. 419, 435 (1982); \textit{Kaiser Aetna v. United States}, 444 U.S. 164, 176, 179--80 (1979)), government-authorized physical invasions---``whether by plane, boat, cable, or beachcomber''---are physical takings (\textit{United States v. Causby}, 328 U.S. 256 (1946); \textit{Nollan v. California Coastal Comm'n}, 483 U.S. 825 (1987)). Here, California's regulation did not restrict how growers farmed; it appropriated for union organizers a formal right ``to literally `take access'\,'' onto private land.

Rejecting the Ninth Circuit's ``24 hours a day, 365 days a year'' test as insupportable, the Court held that duration and intermittency bear only on the \emph{amount} of just compensation, not whether a taking occurred (\textit{United States v. Dow}, 357 U.S. 17, 26 (1958); \textit{Causby}, 328 U.S. at 259). Nor could California evade takings liability because the access right was a ``slight mismatch'' from a formal easement in gross under state law. Finally, the Court distinguished \textit{PruneYard Shopping Center v. Robins}, 447 U.S. 74 (1980), because the shopping center there was generally open to 25,000 public patrons daily, whereas Cedar Point and Fowler's agricultural premises were closed to the public, and distinguished \textit{NLRB v. Babcock \& Wilcox Co.}, 351 U.S. 105 (1956), as a statutory NLRA decision rather than a Takings Clause holding (\href{https://www.law.cornell.edu/supremecourt/text/20-107}{Cornell LII}).

\subsection{Three Limiting Boundaries Preserving Routine Governance}
To dispel fears that its rule would imperil routine government entry onto private land, the Court articulated three critical exceptions (\href{https://www.law.cornell.edu/supremecourt/text/20-107}{Cornell LII}):
\begin{enumerate}
    \item \textbf{Isolated Trespasses vs.\ Takings:} Isolated physical invasions not undertaken pursuant to a granted right of access remain ordinary torts (or subject to causation factors like those in \textit{Arkansas Game and Fish Comm'n v. United States}, 568 U.S. 23 (2012)), not \emph{per se} takings.
    \item \textbf{Background Restrictions of Property Law:} Entries consistent with preexisting common-law privileges---such as nuisance abatement (\textit{Lucas v. South Carolina Coastal Council}, 505 U.S. 1003, 1028--29 (1992)), public or private necessity, and reasonable Fourth Amendment law-enforcement searches---do not take a right the owner ever possessed.
    \item \textbf{Health and Safety Inspection Conditions:} Under the unconstitutional-conditions framework of \textit{Nollan} and \textit{Dolan v. City of Tigard}, 512 U.S. 374 (1994), government may condition permits, licenses, or registrations on reasonable health and safety inspections (e.g., food, pesticide, pharmaceutical, or nuclear safety) bearing an essential nexus and rough proportionality to public risks. Because California's union-access rule fit none of these three categories, it was a \emph{per se} physical taking.
\end{enumerate}

\subsection{Concurring Opinion (Justice Kavanaugh)}

Justice Kavanaugh joined the majority in full, adding that \textit{NLRB v. Babcock \& Wilcox Co.}, 351 U.S. 105 (1956), independently supported the growers because it recognized employers' Fifth Amendment right to exclude union organizers absent a ``necessity'' showing (such as workers living on company property beyond reasonable union communication) (\href{https://www.law.cornell.edu/supremecourt/text/20-107}{Cornell LII}).

\subsection{Dissenting Opinion (Justice Breyer)}

Justice Breyer argued that California's rule merely \emph{regulated} rather than \emph{appropriated} the right to exclude, creating no traditional estate or easement. Under \textit{Loretto}, \textit{PruneYard}, and \textit{Arkansas Game and Fish}, only \emph{permanent} physical occupations trigger \emph{per se} liability, whereas temporary access (here, roughly 4\% of the year) belongs under \textit{Penn Central} balancing. Warning that the majority's new rule threatens myriad workplace, food-safety, environmental, and childcare inspections and replaces a workable framework with indeterminate exceptions, Justice Breyer also noted on remedies that under \textit{Knick v. Township of Scott}, 588 U.S. 180 (2019), California should be able to foreclose injunctive relief by paying compensation (\href{https://www.law.cornell.edu/supremecourt/text/20-107}{Cornell LII}).

\section{Examples: Future Applications}
\begin{enumerate}

\item \textbf{Same-Side (Per Se Taking): Wildlife Surveys on Fenced Timberland.} A State enacts a statute granting certified environmental nonprofits a right to enter private, fenced timberlands (closed to the public) for four hours per week during spring migration to conduct wildlife-habitat surveys, imposing civil penalties on landowners who refuse entry. Under \textit{Cedar Point}, the statute does not merely restrict how owners harvest timber; it appropriates a formal right of physical invasion for third parties. Intermittent duration (four hours weekly) affects only compensation, not liability. Because no background common-law privilege authorizes third-party wildlife surveys and the mandate is untethered to any government permit or benefit, it effects a \emph{per se} physical taking requiring just compensation.

\item \textbf{Same-Side (Per Se Taking): Canvassers in a Gated Community.} A city ordinance requires the owners' association of every private, gated residential community to admit registered political canvassers two weekends a month during election season so they can go door to door, and fines associations that refuse. The streets are private and closed to the public. Under \textit{Cedar Point}, the ordinance does not regulate how owners use their land; it grants third parties a right to physically enter it, appropriating the owners' right to exclude. That the access is limited to a few weekends affects only the amount of compensation. \textit{PruneYard} does not help the city, because \textit{Cedar Point} confined it to property generally open to the public, and no exception applies: canvassing is not a background common-law privilege, and the access is not a condition on a government benefit.

\item \textbf{Opposite-Side (No Per Se Taking): Licensing Inspections and a Zoning Setback.} Now suppose a State requires commercial meatpacking plants, as a condition of obtaining a state food-processing license, to permit state agriculture inspectors to enter processing floors at reasonable business hours to test for bacterial contamination, and separately enacts a zoning ordinance barring new processing plants within 500 feet of residential schools. Here, the outcome flips completely: (1) the inspection requirement falls squarely within \textit{Cedar Point}'s third exception because physical access is conditioned on a regulatory license and bears an essential nexus and rough proportionality to preventing foodborne public-health risks (\textit{Nollan}/\textit{Dolan}); and (2) the 500-foot zoning setback appropriates no right of physical entry at all---it merely restricts the owner's \emph{use} of her own land, triggering deferential \textit{Penn Central} balancing rather than \textit{Cedar Point}'s \emph{per se} rule.

\item \textbf{Opposite-Side (No Per Se Taking): A Highway Crew's One-Time Intrusion.} While repaving a state highway, a state road crew mistakenly parks its trucks and stores gravel on an adjoining farmer's field for two days, damaging part of a crop, and then leaves. No statute or regulation gives the State any right of access to the field. The farmer claims a \emph{per se} physical taking under \textit{Cedar Point}. The claim fails under the Court's first limiting boundary: an isolated physical invasion not undertaken under a granted right of access is ordinarily a trespass, remedied in tort, not an appropriation of the right to exclude. If the intrusion recurred or became part of a deliberate government program, the farmer could argue for a taking, and courts would weigh factors like those in \textit{Arkansas Game and Fish Comm'n v. United States}, 568 U.S. 23 (2012).

\item \textbf{Fence-Sitter (Unsettled): Union Organizers at a Supermarket.} A state statute gives union organizers a right to stand on the private sidewalk and parking lot of any large supermarket, for limited hours, to speak with employees and customers. The store is open to the shopping public every day, but its owner bars all solicitation. The owner argues that \textit{Cedar Point} makes the statute a \emph{per se} taking because it grants third parties a right of access the owner never consented to. The State responds that \textit{Cedar Point} itself distinguished \textit{PruneYard}, which allowed access to a shopping center generally open to the public. Yet organizers are not customers using the property as invited, and the \textit{Cedar Point} majority treated the right to exclude as a fundamental element of property that cannot be balanced away. How far \textit{PruneYard} survives is uncertain.

\end{enumerate}

\section{Critique}
\textbf{Unanswered Remedial Paradox:} As Justice Breyer noted, the Takings Clause does not forbid takings; it requires \emph{just compensation}. Yet the growers sought injunctive relief rather than damages, leaving unresolved how \textit{Cedar Point} interacts with \textit{Knick}'s rule that available compensation remedies foreclose injunctions.

\textbf{Friction Over the Three Exceptions:} Drawing crisp boundaries between ``isolated'' torts and recurring takings, defining which historical common-law privileges count, and distinguishing ``special benefits'' (e.g., permits) from ``basic and familiar uses of property'' (\textit{Horne}) generate ongoing litigation over environmental, labor, and housing regulations.

\textbf{Fragility of \textit{PruneYard}:} Limiting \textit{PruneYard} to businesses ``generally open to the public'' invites future challenges asking whether public-facing retail venues and shopping centers retain a \textit{Cedar Point} right to exclude non-customers.

\section{Key Quotations}
\begin{quote}
``The essential question is not, as the Ninth Circuit seemed to think, whether the government action at issue comes garbed as a regulation (or statute, or ordinance, or miscellaneous decree). It is whether the government has physically taken property for itself or someone else---by whatever means---or has instead restricted a property owner's ability to use his own property.''
\end{quote}
\noindent --- Roberts, C.J. (\href{https://www.law.cornell.edu/supremecourt/text/20-107}{Cornell LII}).

\begin{quote}
``The access regulation appropriates a right to invade the growers' property and therefore constitutes a per se physical taking. \ldots{} There is no reason the law should analyze an abrogation of the right to exclude in one manner if it extends for 365 days, but in an entirely different manner if it lasts for 364.''
\end{quote}
\noindent --- Roberts, C.J. (\href{https://www.law.cornell.edu/supremecourt/text/20-107}{Cornell LII}).

\begin{quote}
``We cannot agree that the right to exclude is an empty formality, subject to modification at the government's pleasure. On the contrary, it is a `fundamental element of the property right,' \textit{Kaiser Aetna}, 444 U.S., at 179--180, that cannot be balanced away.''
\end{quote}
\noindent --- Roberts, C.J. (\href{https://www.law.cornell.edu/supremecourt/text/20-107}{Cornell LII}).

\begin{quote}
``The regulation does not appropriate anything. \ldots{} It gives union organizers the right temporarily to invade a portion of the property owners' land. It thereby limits the landowners' right to exclude certain others. The regulation regulates (but does not appropriate) the owners' right to exclude.''
\end{quote}
\noindent --- Breyer, J., dissenting (\href{https://www.law.cornell.edu/supremecourt/text/20-107}{Cornell LII}).


\end{document}
